\documentclass[dvipdfmx,a4j]{jbook}
\usepackage{wrapfig}
\usepackage{tikz}
\usetikzlibrary{shapes,calc,arrows.meta}
\usepackage{amsmath}
\usepackage{amssymb}
\usepackage{fancyvrb}
\usepackage{tabularx}
%\graphicspath{{figs241004/}}
\usepackage{here}
\usepackage{graphicx}
\usepackage{fancyvrb}
\usepackage{multirow}
%\include{mymacros}
\usepackage{subfig}
%\usepackage[framemethod=tikz]{mdframed}
\newcounter{coffee}
\newcounter{bean}
\usepackage{docmute}
\usepackage{url}
\usepackage{makeidx}
\usepackage{here}
\usepackage{array}
\usepackage{tabularx}
\usepackage{mdframed}
\usepackage{longtable}
\newtheorem{Def}{定義}[chapter]
\newtheorem{Thm}{定理}[chapter]
\newtheorem{Prop}{命題}[chapter]
\newtheorem{Lem}{補題}[chapter]
\newtheorem{Exmp}{例}[chapter]
\newtheorem{Alg}{アルゴリズム}[chapter]

\newtheorem*{repeatedtheorem}{定理 \repeatednumber}
\newtheorem*{repeatedlemma}{補題 \repeatednumber}

\renewcommand{\proofname}{証明}

\newenvironment{repThm}[1]
{\def\repeatednumber{\ref{#1}}\begin{repeatedtheorem}}
{\end{repeatedtheorem}}

\newenvironment{repLem}[1]
{\def\repeatednumber{\ref{#1}}\begin{repeatedlemma}}
{\end{repeatedlemma}}

\newcommand{\tuple}[1]{\langle{#1}\rangle}
\newcommand{\Term}[1]{`\mbox{$\mathtt{#1}$}'}
\newcommand{\NTerm}[1]{\tuple{\mbox{$\mathsf{#1}$}}}
\newcommand{\First}[1]{\mathsf{First}_1({#1})}
\newcommand{\Follow}[1]{\mathsf{Follow}_1({#1})}
\newcommand{\greentext}[1]{{\color[rgb]{0,0.5,0}{#1}}}
\newcommand{\bluetext}[1]{{\color{blue}{#1}}}

\newcommand{\eof}{{\tt \#}}
\newcommand{\ttpl}{{\tt +}}
\newcommand{\ttsb}{{\tt -}}
\newcommand{\ttmu}{{\tt *}}
\newcommand{\ttdv}{{\tt /}}
\newcommand{\ttlp}{{\tt (}}
\newcommand{\ttrp}{{\tt )}}
\newcommand{\tteq}{{\tt =}}
\newcommand{\ttlt}{{\tt <}}
\newcommand{\ttcm}{{\tt ,}}

\newcommand{\ttbrkl}{{\tt [}}
\newcommand{\ttbrkr}{{\tt ]}}

\newcommand{\ttif}{{\tt if}}
\newcommand{\ttfi}{{\tt fi}}
\newcommand{\ttthen}{{\tt then}}
\newcommand{\ttelse}{{\tt else}}
\newcommand{\ttwhile}{{\tt while}}
\newcommand{\ttdo}{{\tt do}}
\newcommand{\ttod}{{\tt od}}

\newcommand{\ttassgn}{{\tt :=}}
\newcommand{\ttsc}{{\tt ;}}

\newcommand{\ttand}{{\tt \&\&}}
\newcommand{\ttor}{{\tt ||}}
\newcommand{\ttneg}{{\tt !}}

\newcommand{\tti}{\texttt{i}}
\newcommand{\ttj}{\texttt{j}}
\newcommand{\ttk}{\texttt{k}}
\newcommand{\ttx}{\texttt{x}}
\newcommand{\tty}{\texttt{y}}
\newcommand{\ttz}{\texttt{z}}

\newcommand{\rppush}[1]{\mathsf{push}\ \texttt{#1}}
\newcommand{\rpassign}{{\tt :=}}
\newcommand{\rpframe}[1]{\draw #1 -- ($#1+(0,1)$) #1 -- ($#1+(1.5,0)$) ($#1+(1.5,0)$) -- ($#1+(1.5,1)$);}
\newcommand{\rpframerow}[2]{\draw ($#1}

\newcommand{\newTemp}[1]{{\tt newTemp(}#1\texttt{)}}
\newcommand{\newTempnull}{{\tt newTemp}}
\newcommand{\newLabel}[1]{{\tt newLabel(}#1\texttt{)}}
\newcommand{\newLabelnull}{{\tt newLabel}}
\newcommand{\generatefo}[4]{{\tt generate(}{#1},=,{#2},{#3},{#4}\texttt{)}}
\newcommand{\generateth}[3]{{\tt generate(}{#1},=,{#2},{#3}\texttt{)}}
\newcommand{\generatetw}[2]{{\tt generate(}{#1},=,{#2}\texttt{)}}
\newcommand{\generateon}[1]{{\tt generate(}{#1}\texttt{)}}
\newcommand{\generateCTL}[3]{{\tt generate(}{#1}.{#2},{#3}\texttt{)}}
\newcommand{\generate}[1]{{\tt generate(}{#1}{\tt )}}
\newcommand{\generatenull}{{\tt generate}}
\newcommand{\generateassgin}[2]{{\tt generate}(`{#1}{\tt =}{#2}')}
\newcommand{\generateif}[2]{{\tt generate}(`{\tt if}\ {#1}\ {\tt goto}\ {#2}')}
\newcommand{\generateifnot}[2]{{\tt generate}(`{\tt if\_not}\ {#1}\ {\tt goto}\ {#2}')}
\newcommand{\generatelabel}[1]{{\tt generate(`label\ }{#1}')}
\newcommand{\generategoto}[1]{{\tt generate(`goto\ }{#1}')}

\newcolumntype{b}{X}
\newcolumntype{s}{>{\hsize=.5\hsize}X}

\newcolumntype{Y}{>{\centering\arraybackslash}X}
\newcommand{\inst}[1]{\begin{tabularx}{.3\textwidth}{|Y|Y|Y|Y|}\hline\texttt{#1}& $b$ & $i$ &$a$\\\hline\end{tabularx}}
\newcommand{\ignored}[1]{\par #1は無視される。}

\newcommand{\bnfnts}[1]{\langle\mathrm{#1}\rangle}
\newcommand{\bnfts}[1]{\texttt{'#1'}}
\newcommand{\bnfarr}{\texttt{::=}}

\newcommand{\domv}[1]{\mathsf{dom}({#1})}
\newcommand{\idomv}[1]{\mathsf{idom}({#1})}
\newcommand{\domgg}{\underline{\gg}}
\newcommand{\domll}{\underline{\ll}}

\newcommand{\loadinst}[2]{\texttt{Load}\>{$#1$}\>{$#2$}}
\newcommand{\oprinst}[3]{\texttt{#1}\>{#2}\>{#3}}

\setcounter{chapter}{7}
% \setcounter{page}{63}
% \newcounter{coffee}
% \newcounter{bean}
%\graphicspath{{figs241108/}}
\begin{document}
\chapter{目的コード生成}

３番地コードとして生成された命令列は３番地コードを解釈する実行系で実行可能である．
３番地コードは計算機のアーキテクチャからは独立した命令であり，変数はメモリ上にあると
して実行される．コンパイラにおける目的コード生成は，中間コードに対応する目的
計算機の命令語を選択する．ここで，現在の計算機の多くのアーキテクチャでは，メモリにある
データを演算器のレジスタにロードして演算を行い，レジスタに得られた結果をメモリに
書き戻す動作をする．演算器のレジスタとメモリの間のデータの交換速度は一般に計算機の
演算に比べて遅いため，プログラムの解析によってレジスタとメモリの間のデータ交換が
できるだけ少なくなるような目的コードを生成することが望ましい．中間コードの変数の数に
比べてレジスタの数は非常に少ないため，実行プログラムにおいてその時点で頻繁に演算に
用いられる変数をレジスタに置くようなコードを生成することによって実行速度の大幅な
改善が可能になる．変数とレジスタの対応を決定することをレジスタ割当という．

一般にプログラムに対して最適なレジスタ割当を行うことは難しく，３番地コードのフローグラフが
与えられている時に割当アルゴリズムがNP完全な問題である\cite{DBLP:journals/cl/ChaitinACCHM81}ため，
目的コード生成での
レジスタ割当は最適性を保証することは困難である．従って，目的コード生成では，部分的な
プログラムに対して必要なレジスタ数を見積もり，近似的に効率のよい割当手法を与える，
３番地コードやSSA形式から変数のデータフローからライフサイクルを解析して，
変数からレジスタへの割当を決定する．

目的コードの命令の選択については目的マシンのアーキテクチャに応じてできるだけサイクル数の
少ない命令（軽い命令）を選択するなどコストを見積もるなどの手法がある．しかし，
実際の計算機ではハードウェアによる並列的な実行などの効果や実行ユニットの数などに依存する
ため，実際的な計算機に対するコード生成の系統的なモデル化は教科書的内容の範囲を越えるため，
本書では扱わない．

\section{算術式のコード生成}

ここでは，算術式に対するコード生成において必要なレジスタの数を見積もる手法を示す．
簡単化のため，以下の範囲に限定する．
\begin{itemize}
\item 演算子は四則演算($\ttpl$, $\ttsb$, $\ttmu$, $\ttdv$)に限定する．
\item 命令はメモリ転送命令($\texttt{Load}\ r,e$)と演算命令($\theta\ r,e$)とする．\\
メモリ転送命令$\texttt{Load}\ r,e$は，$e$で示されるメモリの内容をレジスタ$r$に転送する．
演算命令$\theta\ r,e$は，$e$で示されるメモリの内容をフェッチして$r$と$\theta\in\{\ttpl,\ttsb,\ttmu,\ttdv\}$
で定められた演算を行い，結果を$r$に格納する．
\end{itemize}

例えば，算術式$A\ttpl B\ttsb C$のコードは以下のようになる．


\begin{minipage}{.3\linewidth}
\begin{tabbing}
\quad\qquad\=\qquad\=\\\kill
\loadinst{r1}{A}\\
\oprinst{+}{r1}{B}\\
\oprinst{-}{r1}{C}\\
\end{tabbing}
\end{minipage}

以上の制約のもとでは，レジスタは１個で計算が可能である．
これに対して，$A\ttsb\ttlp B\ttpl C\ttrp$の場合は以下のような命令系列になる．

\begin{minipage}{.3\linewidth}
    \begin{tabbing}
    \quad\qquad\=\qquad\=\\\kill
    \loadinst{r2}{B}\\
    \oprinst{+}{r1}{C}\\
    \loadinst{r1}{A}\\
    \oprinst{-}{r1}{r2}\\
    \end{tabbing}
\end{minipage}

\section{グラフ彩色によるレジスタ割り当て}
\subsection{SSAに対するレジスタ割り当て}
\end{document}
